\textbf{Определение.}

Расширение $K \supset F$ называется \textit{нормальным}, если
если вместе с каждым элементом поле $K$ содержит и все
сопряжённые к нему над $F$.

По теореме об эквивалентных условиях
расширение $K \supset F$ нормально тогда и только тогда,
когда $K$ -- поле разложения многочлена $f(x) \in F[x]$.

а) Пусть $K \supset L \supset F$ -- башня расширений полей, 
$K \supset F$ -- нормальное расширение. 
Тогда $K \supset L$ -- нормальное расширение. 

Если $K \supset F$ -- нормальное расширение, то
$K$ является полем разложения $f(x) \in F[x] \subset L[x]$, а значит
тем более является полем разложения $f(x) \in L[x]$

б) Приведите пример, когда в этой ситуации $L \supset F$ 
не является нормальным.

$\Q(\sqrt[3]{2}, i) \supset \Q(\sqrt[3]{2}) \supset \Q$

$\Q(\sqrt[3]{2}, i) \supset \Q(\sqrt[3]{2})$ -- нормальное расширение
(является полем разложения $x^2 + 1$)

$\Q(\sqrt[3]{2}, i) \supset \Q$ нормально, так как является полем
разложения $x^3 - 2$

$\Q(\sqrt[3]{2}) \supset \Q$ не является нормальным расширением,
так как $\sqrt[3]{2} \omega$ является сопряжённым к $\sqrt[3]{2}$
над $\Q$, но не лежит в $\Q(\sqrt[3]{2})$.

